\documentclass[10pt,a4paper]{amsart}
\usepackage{amsmath,amssymb,amsthm,graphicx}
\usepackage{verbatim}
\newenvironment{reference}{\comment}{\endcomment}
%\newenvironment{reference}{}{}
\newenvironment{slogan}{\comment}{\endcomment}
\newenvironment{history}{\comment}{\endcomment}

% For commutative diagrams we use Xy-pic
\usepackage[all]{xy}

% We use 2cell for 2-commutative diagrams.
\xyoption{2cell}
% Theorem environments.
%
\theoremstyle{plain}
\newtheorem{theorem}[subsection]{Theorem}
\newtheorem{proposition}[subsection]{Proposition}
\newtheorem{lemma}[subsection]{Lemma}

\theoremstyle{definition}
\newtheorem{definition}[subsection]{Definition}
\newtheorem{example}[subsection]{Example}
\newtheorem{exercise}[subsection]{Exercise}
\newtheorem{situation}[subsection]{Situation}

\theoremstyle{remark}
\newtheorem{remark}[subsection]{Remark}
\newtheorem{remarks}[subsection]{Remarks}

\numberwithin{equation}{subsection}

% Macros
%
\def\lim{\mathop{\mathrm{lim}}\nolimits}
\def\colim{\mathop{\mathrm{colim}}\nolimits}
\def\Spec{\mathop{\mathrm{Spec}}}
\def\Hom{\mathop{\mathrm{Hom}}\nolimits}
\def\Ext{\mathop{\mathrm{Ext}}\nolimits}
\def\SheafHom{\mathop{\mathcal{H}\!\mathit{om}}\nolimits}
\def\SheafExt{\mathop{\mathcal{E}\!\mathit{xt}}\nolimits}
\def\Sch{\mathit{Sch}}
\def\Mor{\mathop{\mathrm{Mor}}\nolimits}
\def\Ob{\mathop{\mathrm{Ob}}\nolimits}
\def\Sh{\mathop{\mathit{Sh}}\nolimits}
\def\NL{\mathop{N\!L}\nolimits}
\def\CH{\mathop{\mathrm{CH}}\nolimits}
\def\proetale{{pro\text{-}\acute{e}tale}}
\def\etale{{\acute{e}tale}}
\def\QCoh{\mathit{QCoh}}
\def\Ker{\mathop{\mathrm{Ker}}}
\def\Im{\mathop{\mathrm{Im}}}
\def\Coker{\mathop{\mathrm{Coker}}}
\def\Coim{\mathop{\mathrm{Coim}}}

% Boxtimes
%
\DeclareMathSymbol{\boxtimes}{\mathbin}{AMSa}{"02}

%
% Macros for moduli stacks/spaces
%
\def\QCohstack{\mathcal{QC}\!\mathit{oh}}
\def\Cohstack{\mathcal{C}\!\mathit{oh}}
\def\Spacesstack{\mathcal{S}\!\mathit{paces}}
\def\Quotfunctor{\mathrm{Quot}}
\def\Hilbfunctor{\mathrm{Hilb}}
\def\Curvesstack{\mathcal{C}\!\mathit{urves}}
\def\Polarizedstack{\mathcal{P}\!\mathit{olarized}}
\def\Complexesstack{\mathcal{C}\!\mathit{omplexes}}
% \Pic is the operator that assigns to X its picard group, usage \Pic(X)
% \Picardstack_{X/B} denotes the Picard stack of X over B
% \Picardfunctor_{X/B} denotes the Picard functor of X over B
\def\Pic{\mathop{\mathrm{Pic}}\nolimits}
\def\Picardstack{\mathcal{P}\!\mathit{ic}}
\def\Picardfunctor{\mathrm{Pic}}
\def\Deformationcategory{\mathcal{D}\!\mathit{ef}}
\title{A projective scheme and a finite type quasi-coherent module on P\textasciicircum{}1 with non-finitely generated H\textasciicircum{}0}
\author{The Stacks Project Authors}
\date{stacks-project:examples.tex:local-snapshot-20261012; GNU FDL 1.2 or later}
\begin{document}
\maketitle
\noindent\textit{Editorial note.} All delivered mathematical text is the author's own, copied from the pinned Stacks Project chapter file examples.tex; no statement or proof step was rewritten, reordered or removed.\par\medskip
\noindent\textit{Editorial note.} The counted proof is the author's complete argument: The counted proof adds no mathematics of its own; the bidegree table at 5299-5313 is the pivot of the kernel computation.\par\medskip
\noindent\textit{Editorial note.} The label lemma-nonfinite-h0 names the counted statement, and the proof environment immediately following it in the source is delivered verbatim as its proof.\par\medskip
\noindent\textit{Editorial note.} The counted statement is carried by the author's own lemma environment in the chapter source.\par\medskip
\noindent\textit{Allowed standard prerequisites (editorial).}\begin{itemize}
\item Context span 5280-5285: R any non-Noetherian ring, I an ideal that is not finitely generated, S = R[x, y]/xyI graded with x, y of degree 1; X = Proj(S) is a projective scheme over R; the claim is that H\textasciicircum{}0(X, O\_X) is not a finitely generated R-module
\item Context span 5287-5325: decomposing S\_x, S\_y, S\_\{xy\} by bidegree gives summands R'\_\{a,b\} (R for b = 0, R/I for b != 0), R''\_\{a,b\} (R for a = 0, R/I for a != 0), and R/I in every bidegree; H\textasciicircum{}0(X, O\_X) is the kernel of (S\_x)\_0 x (S\_y)\_0 -> (S\_\{xy\})\_0, (F, G) -> F - G, which by the explicit descriptions is the submodule M of R\textasciicircum{}2 of pairs with F - G in I; the surjection M -> I given by (F, G) -> F - G shows M is not finitely generated
\item Context span 5327-5332: the surjection R[x, y] -> S gives a closed immersion i: X -> P\textasciicircum{}1\_R; F = i\_*O\_X is a finite type quasi-coherent module on P\textasciicircum{}1\_R whose H\textasciicircum{}0 equals the non-finitely generated module M
\end{itemize}\medskip
Let $R$ be any non-Noetherian ring and $I$ an ideal of $R$
that is not finitely generated. Let $S = R[x, y]/xyI$ viewed
as a graded ring with $x$ and $y$ having degree $1$.
Then $X = \text{Proj}(S)$ is a projective scheme over $R$.
We claim that $H^0(X, \mathcal{O}_X)$ is not
a finitely generated $R$-module.

\medskip\noindent
Decomposing the rings $S_x$, $S_y$, and $S_{xy}$ by bidegree,
we can write
$$
S_x = \bigoplus_{\substack{(a, b) \in \mathbf{Z}^2 \\ b \geq 0}} R'_{a, b},
\quad\quad
S_y = \bigoplus_{\substack{(a, b) \in \mathbf{Z}^2 \\ a \geq 0}} R''_{a, b},
\quad\quad
S_{xy} = \bigoplus_{\substack{(a, b) \in \mathbf{Z}^2}} R/I,
$$
where the $(a, b)$-summands sit in degree $a + b$ and we have
$$
R'_{a, b} =
\left\{
\begin{matrix}
R & b = 0, \\
R/I & b \ne 0,
\end{matrix}
\right.
\quad\quad
R''_{a, b} =
\left\{
\begin{matrix}
R & a = 0, \\
R/I & a \ne 0.
\end{matrix}
\right.
$$
The $R$-module $H^0(X, \mathcal{O}_X)$ can be computed as the kernel of the map
$$
(S_x)_0 \times (S_y)_0 \to (S_{xy})_0
$$
defined by $(F, G) \mapsto F - G$. Here the notation $A_0$ refers to the
degree $0$ component of a $\mathbf{Z}$-graded ring $A$. By our explicit
description of $S_x, S_y, S_{xy}$ above, this kernel is isomorphic
to the submodule $M \subset R^2$ consisting of those pairs $(F, G)$
with $F - G \in I$. Since the map $M \to I$ of $R$-modules
given by $(F, G) \mapsto F - G$ is surjective,
$M$ is not a finitely generated $R$-module.

\medskip\noindent
The surjection $R[x, y] \to S$ corresponds to a closed
immersion $i : X \to \mathbf{P}^1_R$. Then $\mathcal{F} = i_*\mathcal{O}_X$
is a finite type quasi-coherent $\mathcal{O}_{\mathbf{P}^1_R}$-module
whose $H^0$ is equal to the non-finitely generated module $M$
we computed above.

\begin{lemma}
\label{lemma-nonfinite-h0}
Non-finite $H^0$.
\begin{enumerate}
\item There exists a ring $R$ and a projective scheme $X$
over $R$ such that $H^0(X, \mathcal{O}_X)$ is not
a finite $R$-module.
\item There exists a ring $R$ and a finite type quasi-coherent module
$\mathcal{F}$ on $\mathbf{P}^1_R$ such that
$H^0(\mathbf{P}^1_R, \mathcal{F})$ is not a finite $R$-module.
\end{enumerate}
\end{lemma}

\begin{proof}
See discussion above.
\end{proof}

\end{document}
